\documentclass{extreport}

\usepackage[14pt]{extsizes}
\usepackage[T2A]{fontenc}
\usepackage[utf8]{inputenc}
\usepackage[english,ukrainian]{babel}

\usepackage[a4paper, top=10mm, bottom=15mm, left=20mm, right=15mm]{geometry}
\usepackage{amsmath,amsfonts,amssymb,amsthm,mathtools}
\usepackage{listings}
\usepackage{graphicx}
\usepackage{enumitem}
\usepackage{verbatim}
\usepackage{listings}
\usepackage{xcolor}
\usepackage{textgreek}
\usepackage{diagbox}
\usepackage{pgfplots}
\pgfplotsset{compat = 1.16}


\lstdefinestyle{output_txt}{
    basicstyle=\ttfamily\footnotesize,
    breakatwhitespace=false,         
    breaklines=true,                 
    captionpos=b,                    
    keepspaces=true,                                    
    numbersep=5pt,                  
    showspaces=false,                
    showstringspaces=false,
    showtabs=false,                  
    tabsize=2
}
\definecolor{codegreen}{rgb}{0,0.6,0}
\definecolor{codegray}{rgb}{0.5,0.5,0.5}
\definecolor{codepurple}{rgb}{0.58,0,0.82}
\lstdefinestyle{python_code}{ 
    commentstyle=\color{codegreen},
    keywordstyle=\color{magenta},
    numberstyle=\tiny\color{codegray},
    stringstyle=\color{codepurple},
    basicstyle=\ttfamily\footnotesize,
    breakatwhitespace=false,         
    breaklines=true,                 
    captionpos=b,                    
    keepspaces=true,                            
    numbersep=5pt,                  
    showspaces=false,                
    showstringspaces=false,
    showtabs=false,                  
    tabsize=4
}

\setlist[enumerate]{nosep}
\graphicspath{{pics/}}
\DeclareGraphicsExtensions{.png}

\begin{document}
\begin{titlepage}
    \thispagestyle{empty}
    \begin{center}
        \includegraphics[width = \textwidth]{kpi}
        Міністерство освіти і науки України\\
        Національний технічний університет України\\
        <<Київський політехнічний інститут ім. І. Сікорського>>\\
        Інститут прикладного системного аналізу
    \end{center}
    \vspace{40mm}
    \begin{center}
        \textbf{Лабораторна робота} \\
        з курсу <<Методи оптимізації>> \\
        з теми <<Методи спряжених градієнтів>
    \end{center}
    \vspace{20mm}
    \begin{flushleft}
        Виконали студенти 3 курсу групи КА-81 \\
        Галганов Олексій \\
        Єрко Андрій \\
        Фордуй Нікіта
    \end{flushleft}
    \begin{flushright}
        Перевірили \\
        Спекторський Ігор Якович \\
        Яковлева Алла Петрівна
    \end{flushright}
    \vspace{30mm}
    \begin{center}
        \textbf{Київ 2021}
    \end{center}
\end{titlepage}

\begin{center}
    \textbf{Варіант 1}
\end{center}
\noindent\textbf{Завдання.} Скласти програму для мінімізації цільової функції методом спряжених градієнтів.

\emph{Цільова функція:}
$f(x,y) = x^2 + 18y^2 + 0.01xy + x - y$

\noindent\textbf{Результати роботи.}
Цільова функція є квадратичною: $$f(x, y) = \frac{1}{2} \left<\begin{pmatrix}
    2 & 0.01 \\ 0.01 & 36
\end{pmatrix} \begin{pmatrix}
    x \\ y
\end{pmatrix}, \begin{pmatrix}
    x \\y
\end{pmatrix} \right> + 
\left<\begin{pmatrix}
    1 \\ -1
\end{pmatrix}, \begin{pmatrix}
    x \\y
\end{pmatrix} \right>$$
Мінімум --- у точці $\begin{pmatrix}
    -0.50014 \\
    0.0279167
\end{pmatrix}$.
Метод спряжених напрямків має мінімізувати її за не більше ніж 2 кроки. Дійсно, при запуску методу з випадкових початкових точок
отримали точку мінімуму за два кроки.
\begin{center}
    \begin{tabular}{c c}
        \includegraphics[scale = 0.5]{quadratic_random_1.png} &
        \includegraphics[scale = 0.5]{quadratic_random_2.png} \\
        \includegraphics[scale = 0.5]{quadratic_random_3.png} &
        \includegraphics[scale = 0.5]{quadratic_random_4.png} \\
        \includegraphics[scale = 0.5]{quadratic_random_5.png} &
        \includegraphics[scale = 0.5]{quadratic_random_6.png} \\
    \end{tabular}
\end{center}

Ми також вирішили перевірити цей метод на функції Розенброка
$f(x, y) = (y-x^2)^2 + (1-x)^2$, яка вже не є квадратичною. Вона має мінімум у точці $(1, 1)$.
Перевірку зробили з чотирьох початкових точок.
\begin{center}
    \begin{tabular}{c c}
        \includegraphics[scale = 0.5]{rosenbrok_0.png} &
        \includegraphics[scale = 0.5]{rosenbrok_1.png} \\
        \includegraphics[scale = 0.5]{rosenbrok_10.png} &
        \includegraphics[scale = 0.5]{rosenbrok_100.png} \\
    \end{tabular}
\end{center}

\noindent\textbf{Блок-схема алгоритму (для неквадратичної функції).}\\
\begin{center}
    \includegraphics[width=0.9\textwidth]{diagram.png}
\end{center}

\noindent\textbf{Лістинг.}
Текст програми було розділено на \texttt{Optimizer.py} з реалізацією
власне методу спряжених градієнтів для випадку квадратичних та неквадратичних функцій,
а також методу дихотомії для розв'язку задачі одновимірної оптимізації
та \texttt{lab4.py}, де викликаються необхідні функції та зберігаються результати.

\noindent\texttt{Optimizer.py}
\lstinputlisting[language = Python, style = python_code]{../code/Optimizer.py}

\noindent\texttt{lab4.py}
\lstinputlisting[language = Python, style = python_code]{../code/lab4.py}

\noindent\textbf{Висновки.} При виконанні даної роботи ми дослідили
застосування методу спряжених градієнтів до мінімізації квадратичної функції та функції Розенброка.
Для квадратичної функції, як і очікувалось, метод збігався до точки мінімуму не більше ніж 2 кроки.
Для функції Розенброка використовувався варіант методу з оновленням кроків, а для розв'язку задачі
одновимірної оптимізації --- метод дихотомії. В цьому випадку метод спряжених градієнтів виявився
суттєво швидшим за градієнтий метод та дещо повільнішим за метод Ньютона. При цьому слід зауважити, що
на відміну від метода Ньютона він потребує значно менше обрахунків, а тому є більш бажаним для мінізації подібних функцій.

\end{document}